\documentclass[11pt,a4paper]{article}
\usepackage[margin=2.4cm]{geometry}
\usepackage{booktabs}
\usepackage{longtable}
\usepackage{amsmath}
\usepackage{xcolor}
\usepackage{enumitem}
\usepackage{fancyhdr}
\usepackage{float}
\usepackage[colorlinks=true,linkcolor=black,urlcolor=blue]{hyperref}

\definecolor{warn}{RGB}{176,0,32}
\definecolor{ok}{RGB}{0,110,60}
\newcommand{\gap}[1]{\textcolor{warn}{\textbf{#1}}}

\pagestyle{fancy}
\fancyhf{}
\rhead{\small Waste-detection dataset plan}
\lhead{\small\thepage}
\setlength{\headheight}{14pt}

\title{\textbf{Building a Training Set for Drone-and-Buggy Litter Detection}\\[2mm]
\large Gap analysis of the current image set and a capture plan}
\author{}
\date{\today}

\begin{document}
\maketitle
\thispagestyle{fancy}

\section*{Summary}

The target system is a three-stage pipeline: a 4K drone scans from high altitude,
descends for a low-altitude pass, and a ground buggy with a camera and arm
retrieves the objects. Each stage sees objects at a completely different scale,
from a completely different angle, and the current image set supports
\emph{none of them}.

Three findings drive everything below.

\begin{enumerate}[leftmargin=*]
\item \gap{There is no aerial imagery at all.} All 244 images are handheld
iPhone photographs taken from roughly waist height. The directory scaffold for
aerial, occlusion, lighting and weather variation exists but contains only 27
\texttt{.gitkeep} placeholder files and zero images.

\item \gap{Detecting cigarette butts from 20\,m is not possible with the Neo 2.}
This is a physical limit, not a model-quality problem (Section~\ref{sec:gsd}).
The Neo 2's 16.5\,mm-equivalent lens is very wide, so at 20\,m a butt spans
about \emph{two} pixels, and even a plastic bottle is marginal. Butts require
flying at 1.5--2\,m; the survey pass should drop from 20\,m to 10--12\,m. No
amount of training changes these numbers.

\item \gap{The set is too small and too unbalanced to train on.} After
classification, 107 of 244 images contain no identifiable object, and 13 of the
23 classes have five examples or fewer. Useful object detection needs on the
order of $10^3$ instances per class, not $10^0$.
\end{enumerate}

The good news is that the hard part --- knowing what the target objects look like
in situ, and a working manual polygon annotation workflow --- already exists and
should be preserved.

\section{What the dataset currently is}

\begin{table}[H]
\centering
\begin{tabular}{lr}
\toprule
\textbf{Property} & \textbf{Value} \\
\midrule
Images & 244 \\
Capture device & iPhone SE (3rd gen), 242/244 \\
Resolution & $4032\times3024$ (12\,MP), 240/244 \\
Images with GPS EXIF & 242 \\
Viewpoint & Handheld, ground level, near-nadir \\
Altitudes represented & One ($\approx$1.2\,m handheld) \\
Manually annotated files & 100 \\
Manual polygons & 331 \\
--- of which are background (rock, grass, road, sky) & 205 \\
--- of which are waste objects & 126 \\
Classes after auto-classification & 23 (+ \texttt{dontknow}) \\
\bottomrule
\end{tabular}
\caption{Current state. Every image comes from a single camera at a single
height, which means the model cannot learn scale or viewpoint invariance from
this data.}
\end{table}

\noindent Class distribution is severely long-tailed:

\begin{center}
\begin{tabular}{lr@{\hskip 1.2cm}lr}
\toprule
\textbf{Class} & \textbf{n} & \textbf{Class} & \textbf{n} \\
\midrule
\texttt{dontknow} (no object found) & 107 & \texttt{glove} & 4 \\
\texttt{cigarette\_butt} & 39 & \texttt{candy\_wrapper} & 4 \\
\texttt{plastic\_bottle} & 17 & \texttt{can} & 4 \\
\texttt{plastic\_bag} & 14 & \texttt{pine\_cone} & 3 \\
\texttt{plastic\_wrapper} & 7 & \texttt{chair} & 3 \\
\texttt{textile}, \texttt{plastic\_piece}, \texttt{paper} & 6 each & \texttt{bottle\_cap} & 3 \\
\texttt{straw}, \texttt{container}, \texttt{cig\_pack} & 5 each & 6 further classes & 1 each \\
\bottomrule
\end{tabular}
\end{center}

\section{The resolution limit sets the mission profile}
\label{sec:gsd}

Ground sample distance (GSD) is how much real-world length one pixel covers:
\[
\mathrm{GSD} = \frac{2h\tan(\theta/2)}{W}
\]
for altitude $h$, horizontal field of view $\theta$, and image width $W$ pixels.

\subsection{Numbers for the DJI Neo 2}

The available airframe is a DJI Neo 2: 1/2-inch 12\,MP CMOS, \textbf{16.5\,mm
equivalent focal length}, f/2.2, FOV quoted as $119.8^\circ$, 4K video at
$3840\times2160$ up to 100\,fps, 2-axis (tilt/roll) mechanical gimbal, downward
infrared plus forward LiDAR obstacle sensing.

This lens is \emph{much wider} than a typical drone camera, which is good for
coverage and bad for resolving small objects. Taking the 16.5\,mm equivalent
gives a horizontal FOV of $95^\circ$; reading DJI's quoted $119.8^\circ$ as a
16:9 diagonal gives $113^\circ$. DJI generally quotes FOV before dewarping, so
the truth sits between --- the table uses $95^\circ$ and the text notes the
pessimistic case.

\begin{table}[H]
\centering
\begin{tabular}{rrr rrrr}
\toprule
\textbf{Alt.} & \textbf{Swath} & \textbf{GSD} & \multicolumn{4}{c}{\textbf{Object size in pixels (longest dimension)}} \\
\cmidrule(l){4-7}
(m) & (m) & (mm/px) & Cig.\ butt & Cap & Wrapper & Bottle \\
\midrule
20   & 43.6 & 11.36 & \gap{2.2}  & \gap{2.6}  & \gap{5.3}  & \gap{22} \\
10   & 21.8 & 5.68  & \gap{4.4}  & \gap{5.3}  & \gap{10.6} & 44 \\
5    & 10.9 & 2.84  & \gap{8.8}  & \gap{10.6} & 21         & 88 \\
3    & 6.5  & 1.70  & \gap{14.7} & \gap{17.6} & 35         & 147 \\
2    & 4.4  & 1.14  & 22         & 26         & 53         & 220 \\
1    & 2.2  & 0.57  & 44         & 53         & 106        & 440 \\
0.5  & 1.1  & 0.28  & 88         & 106        & 211        & 880 \\
\bottomrule
\end{tabular}
\caption{DJI Neo 2, 4K, $95^\circ$ horizontal FOV. Red entries fall below
roughly 20\,px, the practical floor at which small-object detectors still work.
A cigarette butt is about $8\times25$\,mm.}
\end{table}

\noindent Inverting this gives the \textbf{maximum altitude at which each object
class is still detectable} (20\,px threshold). The second column is the
pessimistic $113^\circ$ reading --- plan against it:

\begin{center}
\begin{tabular}{lrr}
\toprule
\textbf{Object} & \textbf{Max.\ alt.\ ($95^\circ$)} & \textbf{Max.\ alt.\ ($113^\circ$)} \\
\midrule
Cigarette butt & 2.2\,m & 1.6\,m \\
Bottle cap     & 2.6\,m & 1.9\,m \\
Wrapper        & 5.3\,m & 3.8\,m \\
Plastic bottle & 22\,m  & 16\,m \\
Chair / large debris & 70\,m+ & 50\,m+ \\
\bottomrule
\end{tabular}
\end{center}

\noindent These ceilings are roughly \textbf{30\% lower than for a conventional
24\,mm-equivalent drone camera}, purely because of the wide lens. The practical
consequences are specific:

\begin{itemize}[leftmargin=*,itemsep=2pt]
\item \gap{Cigarette butts require flying at 2\,m or below} --- realistically
1.5\,m to keep margin. That is a low, slow pass demanding good altitude hold;
the Neo 2's downward IR and LiDAR make it feasible, but it is the binding
constraint on the whole mission.
\item \gap{Even plastic bottles are marginal at 20\,m} (16--22\,px). The high
survey pass should be flown at \textbf{10--12\,m, not 20\,m}, if it is expected
to detect anything at all.
\item The wide lens partly compensates: at 2\,m the swath is 4.4\,m, versus
3.0\,m for a 24\,mm lens. With 20\% overlap that is $\approx$2.9\,km of track per
hectare instead of 3.3\,km --- still substantial, so the low pass must be
targeted rather than exhaustive.
\end{itemize}

\subsection*{Revised mission profile}

The proposed profile --- scan at 20\,m, then 1--2\,m, then buggy --- needs one
correction: \textbf{the 20\,m pass cannot be a litter-detection pass}. It should
be re-scoped as a \emph{survey} pass, flown lower than planned:

\begin{center}
\begin{tabular}{llp{7.4cm}}
\toprule
\textbf{Pass} & \textbf{Altitude} & \textbf{Role} \\
\midrule
Survey  & 10--12\,m & Large items (bottles, bags, cans, furniture); terrain and
surface map; produces the priority regions for the next pass. Not a
small-object pass. \\
\addlinespace
Detection & 1.5--2\,m & Primary small-object detection. Slow, targeted at
regions flagged by the survey pass. \\
\addlinespace
Buggy & 0.2--1\,m & Approach, identification, grasp. \\
\bottomrule
\end{tabular}
\end{center}

\noindent Two Neo-2-specific cautions for the low pass:

\begin{itemize}[leftmargin=*,itemsep=2pt]
\item \textbf{Downwash.} The Neo 2 is a very light airframe, and at 1.5--2\,m
its prop wash may physically move cigarette butts, wrappers and dry leaves.
This corrupts both detection and the ground truth of a seeded test plot. Test
it explicitly before building a protocol on top of it: hover at 2\,m over a
seeded plot and compare before/after layouts.
\item \textbf{2-axis gimbal.} Tilt and roll are stabilised; \emph{yaw is not}.
Expect residual rotational shake, worse in wind and worse at low altitude in
ground effect. Fly the low pass in calm conditions and keep speed down.
\end{itemize}

\noindent A useful lever: the 4K/100\,fps mode allows a much shorter effective
exposure per frame. For a fixed ground speed this cuts motion blur
substantially, which matters more here than frame rate itself --- blur is what
destroys a 22-pixel object.

\section{What needs to be captured}

Targets below assume a detector trained per class. Quantities are
\emph{annotated object instances}, not images; one image can carry many.

\subsection{Altitude and viewpoint coverage}

\begin{center}
\begin{tabular}{llll}
\toprule
\textbf{Tier} & \textbf{Height} & \textbf{Angle} & \textbf{Purpose} \\
\midrule
Buggy-low   & 0.2\,m & 0--30$^\circ$ oblique & Final approach, arm alignment \\
Buggy-high  & 1\,m   & nadir + 45$^\circ$    & Buggy search/navigation \\
Drone-low   & 2\,m   & nadir                 & Primary small-object detection \\
Drone-low   & 5\,m   & nadir + 30$^\circ$    & Transition scale \\
Drone-high  & 10\,m  & nadir                 & Large-object detection \\
Drone-high  & 20\,m+ & nadir                 & Survey, terrain prioritisation \\
\bottomrule
\end{tabular}
\end{center}

\noindent \textbf{The single most valuable capture method} is the
\emph{altitude ladder}: place a known object, then photograph the identical
scene at 0.2, 1, 2, 5, 10 and 20\,m without moving anything. This yields exact
cross-scale correspondence, lets you measure empirically where each class stops
being detectable, and supports scale-augmentation training. Repeat per surface
type. Roughly 50 ladders would transform the dataset's usefulness.

Angles matter because the buggy sees objects obliquely while the drone sees them
from directly above --- these look nothing alike. Every class needs both.

\subsection{Occlusion}

Currently unrepresented as a labelled axis. Capture each class at four levels,
and \textbf{record the occlusion fraction in the annotation} so it can be
evaluated separately:

\begin{center}
\begin{tabular}{ll}
\toprule
\textbf{Level} & \textbf{Definition} \\
\midrule
None    & Fully visible against surface \\
Light   & $<25\%$ hidden; grass blades crossing the object \\
Partial & 25--60\% hidden; in grass tufts, under leaf edges \\
Heavy   & $>60\%$ hidden; only a fragment protrudes \\
\bottomrule
\end{tabular}
\end{center}

Occluders to cover explicitly: tall grass, cut grass, dead leaves, loose soil,
gravel, twigs/pine needles, puddles, and \emph{other litter} (overlapping
objects, which the pipeline will meet in dumping hot-spots).

\subsection{Surface, lighting, weather}

Backgrounds already present are gravel, soil and grass. Missing and needed:
asphalt, concrete/paving, sand, wood chip, leaf litter, and short mown lawn.
Surface matters more than usual here because low-contrast pairings are the hard
cases --- a beige butt on gravel, a clear bottle on wet asphalt.

Lighting in the current set is 230/244 bright sunlight. Required additions:
overcast, deep shadow, dappled shade under trees, golden hour with long shadows,
and dawn/dusk. Add wet-surface and post-rain conditions. If night operation is
intended, capture under drone illumination.

\subsection{Negatives}

The 107 \texttt{dontknow} images are an asset, not waste --- they are
\textbf{hard negatives}: real ground with no litter. Keep them as explicitly
labelled empty frames. Add deliberate confusers: pine cones, stones, leaves,
bark, dog waste, twigs, mushrooms, and flowers, all of which resemble litter at
low pixel counts. Target 20--30\% of the final set as negatives.

\subsection{Volume targets}

\begin{center}
\begin{tabular}{lrr}
\toprule
\textbf{Class tier} & \textbf{Instances (minimum)} & \textbf{Instances (comfortable)} \\
\midrule
Priority (cigarette butt, bottle, can, wrapper, bag) & 1{,}000 & 3{,}000 \\
Secondary (cup, cap, straw, glove, textile, paper)   & 500     & 1{,}500 \\
Rare / large (chair, crate, furniture)               & 200     & 500 \\
Hard negatives (empty ground, confusers)             & 2{,}000 & 5{,}000 \\
\bottomrule
\end{tabular}
\end{center}

Each priority-class instance count should be spread across the altitude tiers,
not concentrated at one height. A practical rule: no more than 40\% of a class's
instances from any single altitude tier.

\section{Annotation requirements}

The existing polygon workflow is sound and should continue --- polygons are worth
the extra effort here because they support segmentation, and grasp planning for
the arm needs object outline, not just a box.

Changes needed:

\begin{enumerate}[leftmargin=*,itemsep=2pt]
\item \gap{Fix the label vocabulary.} Current labels mix Danish and English
(\texttt{serviet}, \texttt{krukker}, \texttt{grankogler}) and contain typos and
truncations (\texttt{dirty\_clot}, \texttt{dirt\_r}, \texttt{bottle\_},
\texttt{gra}, \texttt{1}, \texttt{2}). Freeze a controlled vocabulary and
validate against it at annotation time.

\item \gap{Separate background from object classes.} 205 of 331 polygons label
terrain (rock, grass, road, sky). These are useful for a segmentation task but
must not be mixed into the detector's object classes.

\item \textbf{Record per-instance metadata}: altitude, camera angle, occlusion
level, surface, lighting. Without these the evaluation in Section~\ref{sec:test}
cannot be broken down, and you will not know \emph{where} the model fails.

\item \textbf{Add a scale reference} to capture sessions --- a coloured marker of
known size in frame lets you verify GSD and recover real object dimensions,
which the arm needs for grasp sizing.

\item \textbf{Verify a sample by hand.} Spot-checks during this analysis found
manual labels that disagree with the imagery (one \texttt{paper} polygon
contains a drinking straw; \texttt{dirty\_cloth} is most likely automotive
carpet). Budget a second-pass review.
\end{enumerate}

\section{Test videos for offline evaluation}
\label{sec:test}

Video is the right evaluation medium because the deployed system sees motion,
motion blur and repeated views of the same object --- and because a video lets
you measure \emph{tracking consistency}, which single images cannot.

\subsection*{Neo 2 recording settings}

\begin{center}
\begin{tabular}{lp{9.2cm}}
\toprule
\textbf{Setting} & \textbf{Value and reason} \\
\midrule
Resolution & 4K $3840\times2160$. Never 1080p --- halving resolution doubles
the minimum detectable object size. \\
Frame rate & 60\,fps normally; \textbf{100\,fps} for all moving takes. The
shorter exposure is what preserves small objects, not the frame count. \\
Bitrate/codec & Highest available (80\,Mbps). Compression artefacts and
20-pixel objects interact badly. \\
Colour & Standard profile, not D-Log, unless you will apply the identical
transform at inference. \\
Gimbal & Nadir ($-90^\circ$) for survey and sweep takes; $-45^\circ$ for the
oblique set. \\
Stabilisation & Record at least a third of takes with EIS off, so the model
sees the real vibration signature it will meet in deployment. \\
Flight & Calm conditions for low passes. Log altitude, speed and GPS for
every take. \\
\bottomrule
\end{tabular}
\end{center}

\noindent\textbf{Before each take}, place a known object set and photograph a
ground-truth layout map from directly above, so every object in the scene has a
known identity and position. Keep takes to 60--90\,s. Fly the same plot across
all conditions so that comparisons isolate one variable at a time.

\subsection*{Required takes}

\begin{longtable}{p{3.4cm}p{4.1cm}p{7.3cm}}
\toprule
\textbf{Take} & \textbf{Profile} & \textbf{What it measures} \\
\midrule
\endhead
Altitude descent & Hover, descend 15\,m$\rightarrow$0.5\,m over a fixed seeded plot, pausing 5\,s at 15, 10, 5, 3, 2, 1, 0.5\,m & \textbf{Fly this one first.} It gives the empirical detection threshold per class. The altitude at which each object first appears is the single most useful number you can collect, and it validates or corrects Table~2 against your actual airframe. \\
\addlinespace
High survey & 10--12\,m, 3\,m/s, straight transects & Large-object recall and false-positive rate over open ground. Flown at 10--12\,m rather than 20\,m per Section~\ref{sec:gsd}. \\
\addlinespace
Low sweep & 1.5--2\,m, 1\,m/s, parallel tracks with 20\% overlap (4.4\,m swath) & Primary small-object detection performance; also validates that the swath/overlap plan actually achieves full coverage. \\
\addlinespace
Speed sweep & 2\,m at 0.5, 1, 2, 4\,m/s; repeat at 60 and 100\,fps & Where motion blur breaks detection --- sets the operational speed limit, and quantifies what the 100\,fps mode buys you. \\
\addlinespace
Occlusion plot & 2\,m over a plot seeded at all four occlusion levels & Recall as a function of occlusion; the number to quote when scoping expectations. \\
\addlinespace
Surface transition & 2\,m crossing grass$\rightarrow$gravel$\rightarrow$asphalt & Background-shift robustness within a single continuous clip. \\
\addlinespace
Lighting series & Same plot at dawn, noon, overcast, dusk, wet & Illumination sensitivity with all else held constant. \\
\addlinespace
Buggy approach & 1\,m$\rightarrow$0.2\,m driving toward a target & The handoff: does the object stay tracked as it grows and the angle flattens? \\
\addlinespace
Buggy patrol & 0.5\,m, continuous drive through a seeded area & End-to-end ground-stage recall and duplicate-detection behaviour. \\
\addlinespace
Grasp approach & 0.2\,m static, arm entering frame & Whether detection survives occlusion \emph{by the robot's own arm} --- a failure mode that is easy to forget and fatal in practice. \\
\addlinespace
Confuser plot & 2\,m and 0.5\,m over pine cones, stones, leaves only & False-positive rate on a scene containing zero litter. \\
\addlinespace
Downwash check & Hover 60\,s at 2\,m, then 1\,m, over a seeded plot & Whether the aircraft's own prop wash displaces the targets. Not a detection test --- a validity test for every other low take. \\
\addlinespace
Empty control & Each altitude over verified-clean ground & Baseline false-positive rate. Any detection here is an error, which makes this the cheapest take to score. \\
\bottomrule
\end{longtable}

\noindent \textbf{Seeded plots.} Lay out a $5\times5$\,m plot with a known
inventory --- e.g.\ 20 cigarette butts, 5 bottles, 5 cans, 10 wrappers, 3 bags
--- at recorded positions. Because the count is known exactly, recall can be
computed without annotating every frame. Build three plots: sparse (realistic),
dense (hot-spot), and confuser-heavy.

\subsection*{Scoring the videos}

Annotate every $N$-th frame (10--30) rather than all frames, and report
precision/recall broken down by altitude, occlusion level, surface and class ---
aggregate numbers will hide exactly the failures that matter. Track two
video-specific metrics: \textbf{re-detection stability} (fraction of frames in
which an object, once seen, stays detected) and \textbf{time-to-first-detection}
during descent.

Keep every test video strictly out of training. Given GPS is recorded, split
train/test \textbf{by location}, not randomly --- random splits across frames of
the same site inflate scores badly.

\section{Cheap baselines: Google, OpenAI and Hugging Face models}
\label{sec:baselines}

Before training anything, it is worth knowing what off-the-shelf models already
achieve on this data --- both as a performance floor and as an annotation
accelerator. This can be done for a few dollars if the work is organised
correctly, and it is easy to spend a hundred times more than necessary by
organising it badly.

\subsection{What these models are and are not for}

\gap{None of these are the deployed detector.} The drone and buggy need an
onboard model with millisecond latency and no network. Cloud vision models serve
three other purposes:

\begin{enumerate}[leftmargin=*,itemsep=1pt]
\item \textbf{A baseline floor.} If a zero-shot model already finds 80\% of your
bottles, a custom detector must beat that to justify itself.
\item \textbf{Auto-labelling.} Pre-annotation that a human corrects is several
times faster than annotating from scratch --- the main lever on the annotation
cost of a $10^3$-instance dataset.
\item \textbf{Adjudication.} A second opinion on ambiguous frames.
\end{enumerate}

A crucial distinction: \textbf{general vision--language models (GPT, Gemini) are
not detectors}. They describe an image; they do not reliably produce boxes for
20-pixel objects. A trial run of GPT-4o over this image set at low detail
reported ``no waste item'' for 107 of 244 images, including photographs
containing a plainly visible rusty metal bar. \textbf{Open-vocabulary detectors
are the right comparison} --- they take a text prompt and return boxes, which is
exactly the target task.

\subsection{Hugging Face: the free tier that matters}

Models run locally cost nothing per image, which makes them the correct default
for anything involving video, where frame counts run to the tens of thousands.
All of the following run on a recent laptop:

\begin{center}
\begin{tabular}{llp{6.1cm}}
\toprule
\textbf{Model} & \textbf{Task} & \textbf{Use here} \\
\midrule
OWLv2 & Open-vocab detection & Zero-shot boxes from prompts such as
``cigarette butt''. Primary baseline. \\
Grounding DINO & Open-vocab detection & Stronger phrase grounding; cross-check
against OWLv2. \\
YOLO-World & Open-vocab detection & Real-time; closest in spirit to what will
run onboard. \\
SAM 2 & Segmentation & Turns a box into a polygon --- converts cheap box labels
into the polygon masks this project already uses. \\
Florence-2 & Caption + grounding & Small, fast, dense captioning for
bulk triage. \\
Qwen2.5-VL & VLM & Local stand-in for GPT/Gemini-class description. \\
\bottomrule
\end{tabular}
\end{center}

\noindent The highest-value combination is \textbf{Grounding DINO or OWLv2 for
boxes, then SAM 2 to convert each box to a polygon}, giving pre-annotations in
the project's existing format at zero marginal cost. Run this over every frame
of every test video; it is free, so there is no reason to sample.

\subsection{Cost mechanics of the paid APIs}

Measured token counts for this dataset (1024\,px images, 250-token prompt,
JSON output) give the following for all 244 images:

\begin{center}
\begin{tabular}{lrrrr}
\toprule
\textbf{Model} & \textbf{In/img} & \textbf{Out/img} & \textbf{Standard} & \textbf{Batch ($-50\%$)} \\
\midrule
gpt-5.4-mini      & 1136 & 123 & \$0.34 & \textbf{\$0.17} \\
gpt-4o (detail low) & 300 & 101 & \$0.43 & \$0.21 \\
gpt-4o (detail high) & 980 & 99 & \$0.84 & \$0.42 \\
gpt-5.4           & 1136 & 131 & \$1.17 & \$0.59 \\
Gemini Flash      & 1136 & 130 & \$0.54 & \$0.27 \\
gpt-5.6-sol       & 1136 & 246 & \$3.19 & \$1.59 \\
\bottomrule
\end{tabular}
\end{center}

\noindent Six levers, in order of effect:

\begin{enumerate}[leftmargin=*,itemsep=2pt]
\item \textbf{Batch APIs halve the price.} Both OpenAI and Google discount
asynchronous batch jobs by 50\%. Offline evaluation has no latency requirement,
so there is never a reason to pay standard rates for it.

\item \textbf{Google's free tier is genuinely free.} Gemini Flash models are
available at no cost in AI Studio (order 1{,}500 requests/day). That covers the
entire 244-image set daily for nothing. \gap{Caveat: free-tier content may be
used to improve Google's products.} If the dataset is to remain proprietary, use
the paid tier --- at \$0.27 batched, that is a cheap way to keep it private.

\item \textbf{Downscale before upload.} Images are 12\,MP; models resample them
anyway. Sending 1024\,px costs the same as sending 4032\,px in tokens but a
fraction of the bandwidth, and above the model's internal tile limit the extra
pixels are simply discarded.

\item \textbf{Know which models honour \texttt{detail: low}.} Measured on this
data: GPT-4o honours it (300 vs.\ 980 tokens); \textbf{GPT-5.4 ignores it
entirely} (1136 either way); GPT-5.6-sol honours it (445 vs.\ 1136). Requesting
low detail on a model that ignores it saves nothing and loses accuracy on the
ones that do not.

\item \textbf{Cascade.} Run the free local model on everything, a cheap API
model on what it flags, and an expensive model only on disagreements. In
practice the expensive tier sees 5--10\% of frames.

\item \textbf{Cache every result.} Store raw responses keyed by image hash and
never re-run the same evaluation. Re-running a finished job is the most common
way this budget is wasted.
\end{enumerate}

\subsection{Video without paying for video}
\label{sec:videocost}

\gap{Do not send raw video to a paid API.} Gemini bills video at roughly
5{,}800 tokens per second of 720p footage. A single 90\,s clip is about 521{,}000
tokens --- \$0.78 on Flash, \$0.39 batched --- and the full 13-take suite comes to
roughly \$5 per evaluation pass on the cheapest model, more on anything larger.
That is affordable once and ruinous if iterated, and iteration is the entire
point of an offline test set.

The alternative costs almost nothing and gives strictly more control:

\begin{enumerate}[leftmargin=*,itemsep=1pt]
\item Extract frames locally with \texttt{ffmpeg} at 1--2\,fps.
\item Run the local Hugging Face detectors over \emph{every} frame --- free, and
the only way to measure per-frame tracking stability anyway.
\item Draw a \textbf{stratified sample of 200 frames} spanning altitude,
occlusion, surface and lighting, and send only those to the paid APIs.
\end{enumerate}

\noindent At 200 frames the paid comparison costs \$0.14 (gpt-5.4-mini batched),
\$0.27 (Gemini Flash batched) or \$0.48 (gpt-5.4 batched) \emph{per model}.
Comparing all three across the whole test suite is under a dollar --- against
roughly \$5 per pass for native video on one model alone.

Native video input is still worth one deliberate run, for the questions frames
cannot answer: temporal consistency, and whether a model tracks an object
through occlusion. Budget it as a single experiment, not a routine.

\subsection{Suggested comparison protocol}

\begin{center}
\begin{tabular}{llr}
\toprule
\textbf{Stage} & \textbf{Model} & \textbf{Cost} \\
\midrule
All frames, all videos & OWLv2 / Grounding DINO / SAM 2, local & \$0 \\
All 244 stills & Gemini Flash, free tier & \$0 \\
244 stills, private & Gemini Flash batch & \$0.27 \\
200-frame stratified sample & gpt-5.4-mini batch & \$0.14 \\
200-frame sample, second opinion & gpt-5.4 batch & \$0.48 \\
Disagreements only ($\approx$20 frames) & gpt-5.6-sol & $<$\$0.15 \\
One native-video run & Gemini Flash batch, 1 take & \$0.39 \\
\midrule
\textbf{Total per full evaluation pass} & & \textbf{$\approx$\$1.45} \\
\bottomrule
\end{tabular}
\end{center}

\noindent Score every model on the same held-out frames with the same metric,
broken down by altitude and occlusion. The interesting output is not a
leaderboard but a \textbf{map of where each model fails}, which tells you which
conditions the training set most needs to cover.

\section{Suggested order of work}

\begin{enumerate}[leftmargin=*,itemsep=1pt]
\item Freeze the label vocabulary; fix the existing typo/language labels.
\item Run the free local baseline (Section~\ref{sec:baselines}) over the
existing 244 images. It costs nothing, and if an off-the-shelf detector already
handles a class well, that class needs less new data than the others.
\item Fly one altitude-descent test video. It costs an afternoon and tells you
the real per-class detection ceiling before you invest in bulk capture.
\item Decide the mission profile from that measurement (Section~\ref{sec:gsd}
gives the prediction; the flight gives the truth).
\item Bulk-capture the altitude ladders across surfaces and lighting.
\item Annotate with the extended per-instance metadata.
\item Record the full test-video suite; keep it sealed from training.
\end{enumerate}

\vspace{4mm}
\noindent\rule{\linewidth}{0.4pt}\\
{\small Figures in Sections 1--2 are computed from the current repository
contents and from the stated 4K/24\,mm-equivalent camera assumption; substitute
the real lens parameters when the airframe is chosen, as GSD scales linearly
with focal length.}

\end{document}
